%------------------------------------------------
\section{Numerical Results}
\begin{frame}{\hspace{3cm}}
\centering

\vspace{1.5cm}

\begin{beamercolorbox}[sep=20pt,center,rounded=true,shadow=true]{title}
    {\fontsize{20}{24}\selectfont \textbf{Numerical Results}}
\end{beamercolorbox}

\vspace{0.8cm}


\vfill

\end{frame}

%------------------------------------------------

%------------------------------------------------
\begin{frame}{Purpose of Numerical Examples}
\begin{itemize}
    \item Illustrate the practical performance of the semiparametric density estimator.
    \item Nonparametric component constructed using a \textbf{Gaussian kernel}.
    \item Bandwidth $b$ is selected prior to estimating the mixing parameter $\pi$.
    \item Two datasets are analyzed:
    \begin{itemize}
        \item Yearly maximum wind speed data.
        \item Magnesium concentration in milk powder.
    \end{itemize}
\end{itemize}
\end{frame}

%------------------------------------------------
\begin{frame}{Semiparametric Density Estimator}
\begin{block}{Convex Combination}
\begin{equation}
    g(x,\pi) = \pi f(x,\hat{\theta}) + (1-\pi)\hat{f}(x), 
    \qquad 0 \leq \pi \leq 1.
\end{equation}
\end{block}

\begin{itemize}
    \item $f(x,\hat{\theta})$: Parametric density with parameters estimated by maximum likelihood.
    \item $\hat{f}(x)$: Kernel density estimator.
    \item $\pi$: Mixing parameter estimated numerically.
    \item Interpretation:
    \begin{itemize}
        \item $\pi = 1$: Purely parametric estimator.
        \item $\pi = 0$: Purely nonparametric estimator.
        \item $0 < \pi < 1$: Compromise between the two.
    \end{itemize}
\end{itemize}
\end{frame}

%------------------------------------------------
\begin{frame}{Gaussian Kernel Density Estimator}
\begin{block}{Kernel Estimator}
\begin{equation}
    \hat{f}(x) = \frac{1}{n b} \sum_{i=1}^{n} 
    K\!\left(\frac{x - x_i}{b}\right),
    \qquad 
    K(u) = \frac{1}{\sqrt{2\pi}} e^{-u^2/2}.
\end{equation}
\end{block}

\begin{itemize}
    \item $b$: Bandwidth controlling the smoothness of $\hat{f}(x)$.
    \item A reasonable choice of $b$ is assumed before estimating $\pi$.
    \item The semiparametric estimate tends to:
    \begin{itemize}
        \item The nonparametric estimator as $b \to 0$.
        \item The parametric estimator as $b \to \infty$.
    \end{itemize}
\end{itemize}
\end{frame}

%------------------------------------------------
\begin{frame}{Pseudolikelihood for Estimating $\pi$}
\begin{block}{Pseudolikelihood Function}
\begin{equation}
    \tilde{L}(\pi) = 
    \prod_{i=1}^{n}
    \left[
        \pi f(x_i,\hat{\theta}) + (1-\pi)\hat{f}(x_i)
    \right],
\end{equation}
\end{block}

\begin{itemize}
    \item The estimate $\hat{\pi}$ maximizes $\tilde{L}(\pi)$.
    \item Ensures stability and avoids bias from self-influence.
\end{itemize}
\end{frame}

%------------------------------------------------
\begin{frame}{Example 1: Wind Speed Data}
\begin{itemize}
    \item Data: Yearly maximum wind speeds (North direction) from Sheridan, Wyoming.
    \item Sample size: $n = 20$ with observations:
    \[
    70,\,61,\,61,\,60,\,61,\,63,\,61,\,67,\,61,\,62,\,47,\,67,\,61,\,49,\,55,\,65,\,57,\,51,\,47,\,56.
    \]
    \item Parametric model: \textbf{Gumbel distribution} with unknown location and scale.
    \item The pseudolikelihood function $\tilde{L}(\pi)$ attains its maximum at:
    \[
        \hat{\pi} \approx 0.8.
    \]
    \item Interpretation: The parametric model provides a good fit to the data.
\end{itemize}
\end{frame}

\begin{frame}{Maximum of the Pseudolikelihood (b =.7s) for Wind Speed Data. }
    \begin{figure}
    \centering
    \includegraphics[width=0.75\linewidth]{Screenshot 2026-04-14 205716.png}
    \label{fig:placeholder}
\end{figure}
\end{frame}

\begin{frame}{Density Estimates for Wind Speed Data (case: $b = 0.7s$)}
    \begin{columns}[c]
        % Left column: Image
        \column{0.6\textwidth}
        \begin{figure}
            \centering
            \includegraphics[width=\linewidth]{image.png}
        \end{figure}

        % Right column: Description
        \column{0.4\textwidth}
        
        \vspace{0.3cm}
        \begin{itemize}
            \item \textbf{Dotted-and-dashed line:} Nonparametric density estimate
            \item \textbf{Dashed line:} Parametric density estimate
            \item \textbf{Solid line:} Semiparametric density estimate
        \end{itemize}
    \end{columns}
\end{frame}


%------------------------------------------------
\begin{frame}{Example 2: Magnesium in Milk Powder}
\begin{itemize}
    \item Data: 22 inductively coupled plasma spectroscopy measurements (micrograms per gram):
    \[
    1201.1,\,1213.3,\,1247.6,\,1177.1,\,1248.9,\,1224.1,\,1240.5,\,1168.4,\,1204.4,1210.1,\,1206.5,\\\hspace{10 mm}1210.0,\,1185.6,\,1177.7,\,1265.7,\,1186.4,\,1256.7,\,1202.3,
    1213.7,\,1186.3,\,1204.9,\,1196.7.
    \]
    \item Parametric model: \textbf{Normal distribution} with unknown mean and variance.
    \item As b tends to 0 the nonparametric and semiparametric estimates exhibit \textbf{multiple modes}, and this is an indication of difficulty with the measurement process that produced the data.
    \item Interpretation: Indicates possible complexities or periodic trends in the measurement process.
\end{itemize}
\end{frame} 

\begin{frame}{Maximum of the log Pseudolikelihood (b =.7s) for Magnesium in Milk Powder Data.}
    \begin{figure}
    \centering
    \includegraphics[width=0.67\linewidth]{WhatsApp Image 2026-04-16 at 16.09.54.jpeg}
    \label{fig:placeholder}
\end{figure}
\end{frame}

%------------------------------------------------
\begin{frame}{Behavior of the Semiparametric Estimator}
\begin{itemize}
    \item The semiparametric estimator provides a compromise between parametric and nonparametric approaches.
    \item Relationship with bandwidth $b$:
    \begin{itemize}
        \item $\pi = 0 \Rightarrow$ semiparametric estimate equals the nonparametric estimate with bandwidth $b=0.5s$.
        \item $\pi = 1 \Rightarrow$ semiparametric estimate equals the parametric estimate with bandwidth $b=0.9s$.
        \item As $b \to 0$, the estimator tends to the nonparametric density.
        \item As $b \to \infty$, it approaches the parametric density.
    \end{itemize}
    \item This behavior parallels penalized likelihood methods for small and large penalties.
\end{itemize}
\end{frame}
